% はじめに。合本では番号なしの章として \input される。
\clearpage
\null\vspace{6mm}%
\noindent{\color{kkaccent}\rule{\linewidth}{1.4pt}}\par\vspace{4mm}%
\noindent{\kkgt{7}\fontsize{21pt}{27pt}\selectfont はじめに}\par\vspace{4mm}%
\noindent{\color{kksecondary}\rule{\linewidth}{.5pt}}\par
\vspace{6mm}%
\newcommand{\kkintrosec}[1]{\par\vspace{1.2\baselineskip}\noindent{\kkgt{5}\large\color{kkaccent}#1}\par\vspace{.4\baselineskip}}

\begin{chapintro}
この冊子は、メタ解析をこれから始める臨床検査技師のために書きました。前半で介入の効果を統合する方法を、後半で検査の精度を統合する方法を、架空の例題と実際のデータの両方を使って、R で手を動かしながら説明します。
\end{chapintro}

\kkintrosec{この冊子が想定している読者}

想定している読者は、検査の測定原理、精度管理、基準範囲、カットオフについては日常の業務でよく知っているものの、メタ解析の論文はあまり読んだことがない検査技師です。統計の式は追わなくても読めるように書きました。R はほとんど初めてという前提で、\kkchref{A}と\kkchref{5}で基本の操作から説明します。

検査技師がメタ解析に関わる場面は、検査の精度を調べるときだけではありません。検査部で新しい測定項目を導入するかどうかを決めるとき、医師との共同研究でデータの整理と解析を担当するとき、診療ガイドラインの推奨がどの研究をもとにしているかを確かめるときにも、メタ解析を読む力と自分で計算する力が役に立ちます。第 1 部の例題に降圧薬の臨床試験を選んだのは、介入の効果を統合する考え方が、検査の精度を統合する方法を理解するための前提になるからです。

\kkintrosec{この冊子の読み方}

第 1 部（第 1〜8 章）は、介入の効果を統合するメタ解析を、レビューの工程の順に説明します。問いを立て、文献を探して選び、データを取り出し、統合し、ばらつきを調べ、結果の確かさを評価し、報告するところまでを一通り扱います。第 1 部だけを読んでも、介入のメタ解析を最初から最後まで自分で実行できるように書きました。

第 2 部（第 9〜15 章）は、検査の精度を統合するメタ解析を説明します。第 2 部は第 1 部の第 1〜5 章を読んでいることを前提にしています。検査のメタ解析は、文献の探し方や選び方は介入の場合とほぼ同じで、統計の部分が大きく違います。第 2 部ではその違いを中心に書きました。

各章の頭には「この章で分かること」という囲みを置き、各章の最後の「この章のまとめ」には、その章の要点と覚えておく判断を数行で書きました。時間がないときは、囲みと最後の節だけを拾い読みしても全体の流れがつかめます。

本文の途中にある「（→ 5.1 節（p. 42））」の形の印は、その内容を説明した節と、その節の始まるページを示します。言葉や考え方が分かりにくくなったときは、印の節に戻って読み直してください。途中で読むのをやめても、そこまでに読んだ章の手順はそのまま使えるように、章ごとに話を閉じてあります。

見出しの前や枠の頭に「発展」の印がついた節と枠は、初めて読むときは飛ばして構いません。印のない部分だけを読めば、メタ解析の論文を読む力と、解析者の作業を手伝うときの手順が身につきます。2 回目に読むとき、または自分で解析するときに、印のある部分へ戻ってください。

\kkintrosec{例題について}

この冊子では、2 種類の例題を並べて使います。

\begin{description}
\item[架空の例題] 第 1 部では、架空の降圧薬 A とプラセボを比べた 10 本のランダム化比較試験を使います。アウトカムは脳卒中の発症です。第 2 部では、発熱のある患者の菌血症をプロカルシトニン（PCT）で診断する精度を調べた、架空の 10 本の研究を使います。
\item[実データの例題] 第 1 部では、BCG ワクチンによる結核予防の 13 本の試験（Colditz 1994）を、第 2 部では、飲酒問題を拾い上げる質問票 AUDIT-C の 14 本の研究（Kriston 2008）を使います。どちらも R のパッケージに収められていて、誰でも同じ結果を再現できます。
\end{description}

架空の例題を使うのは、文献の検索から報告までのすべての段階を、同じ例題で通して見せられるからです。実際の論文では、抄録の選別で除外された文献や、本文から数字を取り出す前の記載はふつう公開されません。架空の例題なら、そうした途中の段階も具体的に示せます。著作権の制約を受けずに抄録や本文の抜粋を載せられる利点もあります。

架空の試験と研究には「Aoki 2009」「Kato 2010」のような名前を付けましたが、実在の論文ではありません。架空のデータの数値は、理解しやすいように作ったものであり、降圧薬や PCT の実際の効果や精度を表しているわけではありません。実在の論文を取り上げるときは、その都度、出典を示します。

\kkintrosec{R とコードについて}

本文に載せた R のコードは、そのまま写して実行すると本文と同じ結果が出るように書きました。コードは淡い色の地に左側の青い線で、R の出力は細い枠で囲んで区別しています。

\begin{rcode}
library(meta)            # パッケージを読み込む
m <- metabin(ev_t, n_t, ev_c, n_c,
             studlab = study, data = bp, sm = "RR")
\end{rcode}

\begin{rout}
                         RR             95% CI     z  p-value
Random effects model 0.7281 (0.6137 to 0.8637) -3.64   0.0003
\end{rout}

R の入手、パッケージのインストール、作業用フォルダの作り方は\kkchref{A}にまとめました。コードの全文も付録 A に載せてあります。

\kkintrosec{言葉の約束}

系統的レビュー（システマティックレビュー）は、あらかじめ決めた方法で関連する研究を漏れなく集め、選び、評価する一連の作業を指します。メタ解析は、集めた研究の結果を統計の方法で一つにまとめる計算を指します。メタ解析は系統的レビューの一部として行うのがふつうで、系統的レビューの中にはメタ解析を行わないもの（研究が少なすぎる、研究どうしが違いすぎる場合）もあります。

「率」と「割合」は区別して使います。割合は全体のうちの一部が占める大きさで、分子が分母に含まれます。率は単位時間あたりの変化で、分母に時間が入ります。そのため、ある時点で疾患をもつ人の割合は「有病割合」と書き、感度や特異度も割合として扱います。陽性の結果のうち本当に疾患をもつ人の割合は「陽性的中割合」と書きます。一般には陽性的中率とも呼ばれます。

統計の用語は、初めて出てきたところで説明し、付録 B の用語集にもまとめました。

\kkintrosec{このあとの章へ}

第 2 部でも、重み付き平均（\kkref{sec:05-3}）、ランダム効果（\kkref{sec:05-4}）、研究間のばらつき（\kkref{sec:06-1}）といった第 1 部の考え方をそのまま使います。第 1 章では、一つの試験だけでは答えが出ない理由と、複数の試験を統合すると何が分かるのかを、降圧薬 A の 10 試験を使って説明します。
